\documentclass[12pt]{article}
\usepackage[utf8]{inputenc}
\usepackage[T1]{fontenc}
\usepackage{hyperref}
\title{Panchpargania Language Resources for K-12 Curriculum Development}
\author{Compiled by Claude Sonnet 5 (claude-sonnet-5) for Andr\'as Kornai}
\date{}
\begin{document}
\maketitle

\section{Introduction}
Panchpargania (ISO 639-3 tdb; Glottocode panc1246) is an Indo-Aryan language of Jharkhand, in the Sadanic group. Some sources treat it as a variety of, or the same as, Kurmali. The 2011 Indian census is reported to have counted about 244,914 Panch Pargania speakers.

\section{Primary Educational Resources}
None found.

\section{Linguistic and Reference Materials}
\begin{itemize}
\item Wikipedia, ``Sadanic languages'': classification and references. \url{https://en.wikipedia.org/wiki/Sadanic_languages}
\item Wikipedia, ``Kurmali language'': mentions the name Panchpargania and gives references. \url{https://wikipedia.com/wiki/Kudmali}
\end{itemize}
No grammar or dictionary specific to Panchpargania was found.

\section{Religious and Community Resources}
\begin{itemize}
\item Joshua Project profile of the Puran people, who are said to speak it around Ranchi. Its figures are unreliable. \url{https://joshuaproject.net/people_groups/17913/IN}
\end{itemize}

\section{Digital and Audio Resources}
None found.

\section{Curriculum Development Context}
Almost nothing specific to Panchpargania exists in the search results. Work on Kurmali and the Sadanic languages, and on the regional-language programmes of Jharkhand, would be the likelier route to teaching material. I did not verify any such programme.
\end{document}